\section{Introduzione}

La presente relazione descrive la progettazione, l'implementazione e la campagna sperimentale del sistema di calcolo del cammino minimo su griglia bidimensionale realizzato per il corso di Algoritmi e Strutture Dati, anno accademico 2024/25.

Il problema affrontato è il calcolo del cammino minimo su una griglia con ostacoli. La difficoltà principale non è geometrica ma computazionale: un'esplorazione ricorsiva ingenua dello spazio dei cammini possibili esplode rapidamente al crescere della griglia, e buona parte del lavoro descritto in questa relazione consiste nel contenere quel costo con strutture dati mirate e regole di potatura.

La figura~\ref{fig:panoramica_sistema} mostra in quattro pannelli, sulla stessa istanza di griglia, il percorso seguito dalla relazione: uno spazio libero senza ostacoli; gli ostacoli generati proceduralmente (Compito 1); la chiusura raggiungibile dall'origine $O$ tramite cammini liberi (in verde/arancio), con la frontiera cerchiata individuata proiettando i raggi cardinali e diagonali (Compito 2); infine il cammino minimo completo verso la destinazione $D$, ottenuto componendo più cammini liberi lungo la sequenza dei landmark (Compito 3). I capitoli seguenti riprendono questi passaggi uno per uno; seguono poi la valutazione delle prestazioni (Compito 4) e le considerazioni finali.

\begin{figure}[H]
\centering
\includegraphics[width=\textwidth]{panoramica_sistema.png}
\caption{Panoramica del sistema sulla stessa griglia: dallo spazio vuoto al cammino minimo risolto.}
\label{fig:panoramica_sistema}
\end{figure}

\subsection{Strumenti e scelte tecnologiche}

Il software è scritto in Python 3.13 e si appoggia a tre librerie principali:
\begin{itemize}
    \item NumPy rappresenta la griglia come array multidimensionale e sfrutta le operazioni vettoriali per le letture e le modifiche dello stato;
    \item Matplotlib genera le visualizzazioni dei percorsi e i grafici dell'analisi delle prestazioni (Compito 4);
    \item Argparse costruisce l'interfaccia a riga di comando.
\end{itemize}

Un approccio a oggetti, con un'istanza dedicata per ogni cella, avrebbe gravato sul gestore della memoria a ogni allocazione; la rappresentazione tramite array numerici lavora invece su dati contigui in memoria.

\subsection{Rappresentazione dello stato della cella}

Ogni cella della griglia è codificata con un numero intero all'interno di una matrice NumPy bidimensionale di tipo \texttt{uint8} (1 byte per cella)~\cite{numpy-dtype}, in base allo schema seguente:

\begin{table}[H]
\centering
\caption{Rappresentazione dello stato della cella a livello logico e fisico}
\begin{tabularx}{\textwidth}{|c|c|X|}
\hline
\textbf{Stato cella} & \textbf{Valore fisico} & \textbf{Significato operativo} \\ \hline
\texttt{Libera} & 0 & Cella attraversabile (navigabile). \\ \hline
\texttt{Ostacolo fisso} & 1 & Ostacolo permanente generato nel Compito 1 o importato da JSON. \\ \hline
\texttt{Ostacolo temporaneo} & $2$ & Cella della chiusura corrente marcata direttamente in memoria durante il ritracciamento ricorsivo; ripristinata a $0$ in fase di risalita. \\ \hline
\end{tabularx}
\label{tab:stato_cella}
\end{table}

Questa codifica numerica rende possibile il \textbf{ritracciamento (backtracking) sul posto}: a ogni chiamata ricorsiva di \texttt{CAMMINOMIN}, il codice sovrascrive direttamente il valore intero della cella visitata nella struttura dati condivisa, invece di clonare la matrice della griglia, operazione che esaurirebbe la memoria sulle griglie più ampie.

Durante la discesa, le celle della chiusura corrente sono marcate con il valore costante $2$; in fase di risalita la funzione le riporta a $0$. Il meccanismo evita di creare e distruggere oggetti a ogni livello dell'albero delle chiamate. Basta una marcatura costante, non legata alla profondità, perché le chiusure annidate sono disgiunte e il ripristino avviene per insieme esplicito; così si evita anche di far eccedere il tipo \texttt{uint8} sulle catene di landmark più lunghe.
